\section*{Summary}

A Dirac material screens the field of the dark matter with its whole dielectric function, $\varepsilon = \varepsilon_{\mathrm{bg}} + \varepsilon_{\mathrm{D}}$ (\cref{sec:16.1}). The published rates screen the Dirac electrons among themselves but not the field of the dark photon, and overstate scattering by $\varepsilon_{\mathrm{bg}}^2$ and absorption by $\varepsilon_{\mathrm{bg}}$ at every mass. For ZrTe$_5$ with the parameters of its measured structure the reach lies near \SI{1.5e-40}{cm^2} at \SI{40}{keV}, 5200 to 6400 times weaker than with the published rule for the standard halo.

Electrons confined to tubes and sheets respond to the momentum along them only, so a halo particle reaches a mode faster than $v_{\mathrm{max}}$ through an oblique momentum transfer, suppressed by the transverse form factor (\cref{sec:16.2}). An electron of graphene that receives more than the work function leaves the sheet with the direction of the momentum transfer, from \SI{1.22}{MeV} upward (\cref{sec:16.3}). The excitons of a molecular crystal follow from the Bethe--Salpeter equation in a molecular basis, and integrated over a window that holds whole manifolds, the crystal responds as a set of independent molecules at every momentum transfer, the Davydov sum rule (\cref{sec:16.4}).

Dilute donors in silicon add a response to first order in their density (\cref{sec:16.5}). The $2p_0$ line of Si:P at \SI{34.11}{meV} reaches $\zeta = \num{2.3e-16}$ at \SI{e15}{cm^{-3}}, twenty times below the continuum, and the first-order form fails at the line center from \SI{4.8e15}{cm^{-3}}. A color center forms after the prompt window: its rate is the recoil rate times the yield $Y_{\mathrm{CC}} = p_{\mathrm{FP}}p_{\mathrm{surv}}p_qf_{\mathrm{opt}}$ (\cref{sec:16.6}). A fluorine recoil in LiF above the simulation threshold of \SI{27}{eV} needs $\mchi \ge \SI{0.186}{GeV}$, and near threshold a recoil makes one vacancy and less than one free electron. Four yields at millikelvin temperatures are open, and each can be measured directly.

Glasses add tunneling two-level systems to the final states (\cref{sec:16.7}). A dark photon is absorbed at a rate proportional to $\zeta^2\tan\delta_0\tanh(m_\phi/2T)/\varepsilon_r$, the dielectric loss of the glass read at the dark-photon mass; on thermal SiO$_2$, three events in a kg-year at \SI{6}{GHz} need $\zeta = \num{2.44e-15}$. Coherent scattering off a random medium leaves the Born approximation once the optical potential exceeds about a tenth of the normal kinetic energy. Two-level systems whose populations froze near \SI{4.8}{K} release heat over hours to weeks after cooling (\cref{sec:16.8}). The density $\bar P$ inferred from heat release in vitreous silica is four times smaller than the one from its specific heat, and the fits to aluminum-coated resonators give $\alpha_{\mathrm{ohm}} = 0.28$ and 0.44, where the small-coupling rate fails. Whether the delayed events of a cooled detector are independent relaxations or avalanches of triggered releases is open, and \cref{sec:19.6} inherits the question.

The uncertainty budget of the chapter is \cref{tab:16-nuisances}, and \cref{tab:16-inputs} lists the material inputs, most of which only a measurement of the sample fixes (\cref{sec:16.9}).
